\chapter{距離空間と測度の準備}
\label{ch:found-preliminaries}

本編で用いる実数・距離空間・測度論の道具をまとめる．
本資料では，使う用語はすべて定義し，
証明なしに用いる標準定理はステートメントを明示したうえで文献を挙げる
（単調収束定理・一致定理・Tonelli の定理・分解定理・
Prokhorov の定理など）．
それ以外の主張にはすべて証明を付ける．
証明を省略した定理については，
例えば伊藤清三『ルベーグ積分入門』，
Dudley \emph{Real Analysis and Probability} (2002)，
Villani \emph{Optimal Transport} (2009) 第1章を参照．

\section{実数の上限・下限}

\begin{definition}{上限と下限}{meas-sup-inf}
  $A\subset\R$ とする．
  $M\in\R$ が $A$ の\textbf{上界}であるとは，
  すべての $a\in A$ で $a\leq M$ が成り立つことをいい，
  上界が存在するとき $A$ は\textbf{上に有界}であるという．
  \begin{itemize}
    \item $A$ が空でなく上に有界であるとき，
      $A$ の上界のうち最小のものがただ一つ存在する
      （実数の連続性公理．本資料では証明せず認める）．
      これを $A$ の\textbf{上限}（supremum）とよび，$\sup A$ と書く．
    \item $A$ が上に有界でないとき $\sup A:=+\infty$，
      $A=\emptyset$ のとき $\sup\emptyset:=-\infty$ と約束する．
  \end{itemize}
  \textbf{下界}・\textbf{下に有界}も同様に定め，
  $A$ が空でなく下に有界であるとき，最大の下界を
  $A$ の\textbf{下限}（infimum）とよび，$\inf A$ と書く．
  $A$ が下に有界でないとき $\inf A:=-\infty$，
  $A=\emptyset$ のとき $\inf\emptyset:=+\infty$ と約束する．
\end{definition}

\begin{definition}{実数列の上極限と下極限}{meas-limsup}
  実数列 $(a_k)$ に対し
  \[
    \limsup_{k\to\infty}a_k
    :=\inf_{n\in\N}\,\sup_{k\geq n}a_k,
    \qquad
    \liminf_{k\to\infty}a_k
    :=\sup_{n\in\N}\,\inf_{k\geq n}a_k
  \]
  と定める（値は $[-\infty,+\infty]$ に取る）．
  つねに $\liminf_ka_k\leq\limsup_ka_k$ であり，
  両者が一致することと $(a_k)$ が（その共通値に）収束することは
  同値である．また，すべての $k$ で $a_k\leq b_k$ ならば
  $\liminf_ka_k\leq\liminf_kb_k$ である（\textbf{$\liminf$ の単調性}）．また
  $k_1<k_2<\cdots$ に対し $(a_{k_j})_{j\in\N}$ を $(a_k)$ の
  \textbf{部分列}という．
  （これらの基本性質は微積分の標準的事実として証明なしに用いる．）
\end{definition}

\begin{remark}{実数についての前提とする事実}{meas-real-facts}
  次はいずれも微積分の標準的事実であり，本資料では証明なしに用いる．
  \begin{itemize}
    \item $\R$ は完備である（Cauchy 列は収束する．連続性公理から従う）．
    \item $p\geq1$ に対し $t\mapsto t^p$ は $[0,\infty)$ 上で
      単調増加かつ凸である．特に
      $\bigl(\frac{a+b}{2}\bigr)^p\leq\frac{a^p+b^p}{2}$（$a,b\geq0$）．
  \end{itemize}
\end{remark}

\section{距離空間}

\begin{definition}{距離空間}{meas-metric-space}
  集合 $M$ 上の関数 $d\colon M\times M\to[0,\infty)$ が
  次の3条件をみたすとき，$d$ を $M$ 上の\textbf{距離}，
  組 $(M,d)$ を\textbf{距離空間}という．
  \begin{enumerate}
    \item （一致の公理）$d(x,y)=0\iff x=y$，
    \item （対称性）$d(x,y)=d(y,x)$，
    \item （三角不等式）$d(x,z)\leq d(x,y)+d(y,z)$．
  \end{enumerate}
\end{definition}

\begin{definition}{開球と開集合}{meas-open-ball}
  距離空間 $(\mathcal{X},d)$，$x\in \mathcal{X}$，$r>0$ に対し
  \[
    B_r(x):=\{y\in \mathcal{X}\mid d(x,y)<r\}
  \]
  を $x$ を中心とする半径 $r$ の\textbf{開球}という．
  $U\subset \mathcal{X}$ が\textbf{開集合}であるとは，
  任意の $x\in U$ に対しある $r>0$ で $B_r(x)\subset U$
  となることをいう．
  補集合が開集合である集合を\textbf{閉集合}という．
\end{definition}

\begin{definition}{連続写像と compact 集合}{meas-continuous-compact}
  距離空間 $\mathcal{X},\mathcal{Y}$ の間の写像 $T\colon \mathcal{X}\to \mathcal{Y}$ が
  \textbf{連続}であるとは，任意の開集合 $V\subset \mathcal{Y}$ に対して
  $T^{-1}(V)$ が $\mathcal{X}$ の開集合であることをいう．

  距離空間 $\mathcal{X}$ の部分集合 $K\subset \mathcal{X}$ が \textbf{compact} であるとは，
  $K$ の任意の開被覆から有限部分被覆が取れることをいう．
  すなわち，開集合族 $(U_i)_{i\in I}$ が
  $K\subset\bigcup_{i\in I}U_i$ をみたすなら，ある有限個
  $i_1,\dots,i_m\in I$ が存在して
  $K\subset U_{i_1}\cup\cdots\cup U_{i_m}$ となる．
  有限個の compact 集合の直積は compact である
  （距離空間の標準的事実として用いる）．
  また，有限個の compact 集合の合併は compact である
  （開被覆を各々に制限して有限部分被覆を合わせればよい）．
\end{definition}


\begin{definition}{可分・完備・Polish 空間}{meas-polish}
  距離空間 $(\mathcal{X},d)$ について：
  \begin{itemize}
    \item $D\subset \mathcal{X}$ が\textbf{稠密}であるとは，
      任意の $x\in \mathcal{X}$ と $r>0$ で $B_r(x)\cap D\neq\emptyset$
      となることをいう．
      可算な稠密部分集合が存在するとき，$\mathcal{X}$ は\textbf{可分}であるという．
    \item 点列 $(x_n)_{n\in\N}$ が \textbf{Cauchy 列}であるとは，
      任意の $\varepsilon>0$ に対しある $N\in\N$ が存在して
      $m,n\geq N$ ならば $d(x_m,x_n)<\varepsilon$ となることをいう．
      すべての Cauchy 列が $\mathcal{X}$ 内の点に収束するとき，
      $\mathcal{X}$ は\textbf{完備}であるという．
    \item 完備かつ可分な距離空間を \textbf{Polish 空間}という．
  \end{itemize}
\end{definition}

\section{測度と可測関数}

$\sigma$-代数・可測空間・測度・可測関数・Lebesgue 積分の定義と，
積分の基本性質（線形性，単調性，可算劣加法性，
零集合上の積分が $0$ であること，
したがって $\mu$-a.e.（定義~\ref{def:meas-ae}）で成り立つ不等式・等式が
積分に遺伝すること，
および非負可測関数の積分が $0$ ならば
その関数が $\mu$-a.e. で $0$ であること）は既知とする．

距離空間 $\mathcal{X}$ の開集合全体を含む最小の $\sigma$-代数を
\textbf{Borel $\sigma$-代数} $\Borel(\mathcal{X})$ とよび，その元を Borel 集合という．
一般に，集合族 $\mathcal{C}$ を含む最小の $\sigma$-代数を
$\mathcal{C}$ が\textbf{生成する} $\sigma$-代数とよび，$\sigma(\mathcal{C})$ と書く．
$\Borel(\mathcal{X})$ 上の測度を Borel 測度といい，
全測度が $\mu(\Omega)=1$ の測度を\textbf{確率測度}，
$\mu(\Omega)<\infty$ の測度を\textbf{有限測度}という
（このとき $(\Omega,\mathcal{A},\mu)$ を有限測度空間という）．
$\mathcal{X}$ 上の Borel 確率測度全体を $\Prob(\mathcal{X})$ と書く．

\begin{remark}{可測関数の基本操作（前提とする事実）}{meas-measurable-ops}
  次はいずれも測度論の標準的な事実であり，本資料では証明なしに用いる．
  \begin{itemize}
    \item 連続関数は Borel 可測である．
      可測関数どうしの和・積・絶対値は可測である．
      可測写像の合成は可測である．
    \item 可測関数列 $(f_n)$ に対し，
      $\sup_n f_n$，$\inf_n f_n$，および各点極限 $\lim_n f_n$
      （存在するとき）は可測である．
    \item 非負可測関数 $f$ に対し，非負単関数
      （可測集合の指示関数の非負線形結合）の列 $(s_n)$ で
      $s_n\uparrow f$（各点）となるものが存在する．
  \end{itemize}
\end{remark}

\begin{theorem}{有限測度の一致定理}{meas-uniqueness}
  集合 $\Omega$ 上の集合族 $\mathcal{C}$ が
  \textbf{$\pi$-系}であるとは，空でなく，
  $A,B\in\mathcal{C}$ ならば $A\cap B\in\mathcal{C}$
  が成り立つことをいう．

  $\mathcal{C}$ を $\Omega$ 上の $\pi$-系とし，
  $\mu,\nu$ を $(\Omega,\sigma(\mathcal{C}))$ 上の有限測度とする．
  $\mu(\Omega)=\nu(\Omega)$ かつ，すべての $A\in\mathcal{C}$ で
  $\mu(A)=\nu(A)$ ならば，$\sigma(\mathcal{C})$ 上で $\mu=\nu$ である．

  （Dynkin の $\pi$-$\lambda$ 定理から従う標準的な結果．証明は省略する．）
\end{theorem}

\section{指示関数と a.e.}

\begin{definition}{指示関数}{meas-indicator}
  集合 $A\subset \mathcal{X}$ の\textbf{指示関数} $\mathbf{1}_A\colon \mathcal{X}\to\{0,1\}$ を
  \[
    \mathbf{1}_A(x):=
    \begin{cases}
      1 & (x\in A),\\
      0 & (x\notin A)
    \end{cases}
  \]
  と定める．$\mu$ が確率測度で $A$ が可測のとき
  $\int_{\mathcal{X}}\mathbf{1}_A\,\d\mu=\mu(A)$ が成り立つ．
\end{definition}

\begin{definition}{ほとんど至るところ（a.e.）}{meas-ae}
  測度空間 $(\Omega,\mathcal{A},\mu)$ において，
  点 $\omega\in\Omega$ に関する性質 $P(\omega)$ が
  $\mu$\textbf{-a.e.}（almost everywhere）で成り立つとは，
  \[
    \mu(\{\omega\in\Omega\mid P(\omega)\text{ が成り立たない}\})=0
  \]
  であることをいう（例外集合が可測であることを含めて要請する）．
\end{definition}

\section{収束定理}

\begin{definition}{各点収束}{meas-pointwise-convergence}
  関数列 $f_n\colon \mathcal{X}\to\R$ が $f\colon \mathcal{X}\to\R$ に\textbf{各点収束}するとは，
  すべての $x\in \mathcal{X}$ に対して
  \[
    \lim_{n\to\infty}f_n(x)=f(x)
  \]
  が成り立つことをいう．
  さらに各 $x$ で $f_1(x)\geq f_2(x)\geq\cdots$ であるとき
  \textbf{単調減少に各点収束する}といい，$f_n\downarrow f$ と書く．
  同様に，各 $x$ で $f_1(x)\leq f_2(x)\leq\cdots$ であるとき
  \textbf{単調増加に各点収束する}といい，$f_n\uparrow f$ と書く．
\end{definition}

\begin{theorem}{単調収束定理}{meas-mct}
  測度空間 $(\Omega,\mathcal{A},\mu)$ 上の非負可測関数列 $(f_n)$ が
  $f\colon\Omega\to[0,+\infty]$ に単調増加に各点収束する
  （$f_n\uparrow f$ と書く．$f(x)=+\infty$ の点では
  $f_n(x)\to+\infty$ の意味に読む）ならば，
  \[
    \lim_{n\to\infty}\int_\Omega f_n\,\d\mu=\int_\Omega f\,\d\mu.
  \]
  （Lebesgue 積分論の基本定理．証明は省略する．）
\end{theorem}

\begin{theorem}{有界収束定理}{meas-bounded-convergence}
  $(\mathcal{X},\Borel(\mathcal{X}),\mu)$ を有限測度空間とする．
  可測関数列 $f_n\colon \mathcal{X}\to\R$ が $f$ に各点収束し，
  ある定数 $M<\infty$ で $\abs{f_n(x)}\leq M$（$\forall n,\,\forall x$）が
  成り立つならば，$f$ も可測で，
  \[
    \lim_{n\to\infty}\int_{\mathcal{X}} f_n\,\d\mu=\int_{\mathcal{X}} f\,\d\mu.
  \]
\end{theorem}

\begin{proof}
  $f$ は可測関数列の各点極限だから可測である
  （注意~\ref{rem:meas-measurable-ops}）．
  \[
    g_n(x) := \sup_{k \geq n} \abs{f_k(x) - f(x)}
  \]
  とおく．$g_n$ は可算個の可測関数の sup だから可測
  （注意~\ref{rem:meas-measurable-ops}）で非負であり，
  $\abs{f_n - f} \leq g_n$ をみたす．
  また $g_n$ は単調減少（$g_1 \geq g_2 \geq \cdots$）であり，
  $f_n \to f$ の各点収束から $g_n(x) \to 0$（すべての $x$）である．
  さらに $\abs{f_k} \leq M$ と $\abs{f} \leq M$ より $g_n \leq 2M$ である．

  $h_n := 2M - g_n$ とおくと，$h_n$ は非負可測かつ単調増加で
  $h_n \uparrow 2M$ である．
  単調収束定理（定理~\ref{thm:meas-mct}）より
  \[
    \int_{\mathcal{X}} h_n \, \d\mu
    \;\uparrow\;
    \int_{\mathcal{X}} 2M \, \d\mu
    = 2M \, \mu(\mathcal{X}).
  \]
  $\mu(\mathcal{X}) < \infty$ だから
  \[
    \int_{\mathcal{X}} g_n \, \d\mu
    = 2M \, \mu(\mathcal{X}) - \int_{\mathcal{X}} h_n \, \d\mu
    \;\to\; 0.
  \]
  $\abs{f_n - f} \leq g_n$ より
  \[
    \left\lvert\int_{\mathcal{X}} f_n \, \d\mu - \int_{\mathcal{X}} f \, \d\mu\right\rvert
    \leq \int_{\mathcal{X}} \abs{f_n - f} \, \d\mu
    \leq \int_{\mathcal{X}} g_n \, \d\mu
    \to 0.  \qedhere
  \]
\end{proof}

\section{確率測度の弱収束と tightness}
\label{sec:meas-weak-tight}

\begin{definition}{有界連続関数の空間}{meas-cb}
  距離空間 $\mathcal{X}$ 上の有界連続関数 $f\colon \mathcal{X}\to\R$ 全体を
  \[
    C_b(\mathcal{X})
    :=
    \{f\colon \mathcal{X}\to\R\mid f\text{ は連続かつ有界}\}
    \quad(\text{有界とは，ある }M<\infty\text{ で }
    \abs{f}\leq M\text{ となること})
  \]
  と書く．
\end{definition}

\begin{definition}{確率測度の弱収束と弱位相（Villani p.~43）}{meas-weak-convergence}
  距離空間 $\mathcal{X}$ 上の Borel 確率測度列 $(\eta_n)$ と
  $\eta\in\Prob(\mathcal{X})$ に対し，
  \[
    \boxed{
      \eta_n\weakto\eta
      \quad\overset{\mathrm{def}}{\Longleftrightarrow}\quad
      \int_{\mathcal{X}}f\,\d\eta_n
      \longrightarrow
      \int_{\mathcal{X}}f\,\d\eta
      \quad(\forall f\in C_b(\mathcal{X}))
    }
  \]
  と定める．このとき $(\eta_n)$ は $\eta$ に\textbf{弱収束}するという．
  ここで $f$ は確率測度を調べる有界連続な\textbf{試験関数}である．

  さらに，すべての $f\in C_b(\mathcal{X})$ に対して
  \[
    \Prob(\mathcal{X})\ni\rho
    \longmapsto
    \int_{\mathcal{X}}f\,\d\rho\in\R
  \]
  が連続となる最も粗い位相を，$\Prob(\mathcal{X})$ 上の
  \textbf{弱位相}という．この位相における列の収束が，上で定めた弱収束である．
\end{definition}

\begin{theorem}{有界連続関数による確率測度の判定}{meas-cb-determines}
  $\mathcal{X}$ を距離空間，$\eta,\rho\in\Prob(\mathcal{X})$ とする．
  すべての $f\in C_b(\mathcal{X})$ に対して
  \[
    \int_{\mathcal{X}} f\,\d\eta=\int_{\mathcal{X}} f\,\d\rho
  \]
  が成り立つならば，$\eta=\rho$ である．
\end{theorem}

\begin{proof}
  閉集合 $F\subset \mathcal{X}$ を任意に取る．$F=\emptyset$ のときは自明なので，
  $F\neq\emptyset$ とする．各 $n\in\N$ に対し
  \[
    f_n(x):=\max\{1-n\,d(x,F),0\},
    \qquad
    d(x,F):=\inf_{y\in F}d(x,y)
  \]
  とおく．距離の三角不等式より
  $\abs{d(x,F)-d(y,F)}\leq d(x,y)$ だから，$f_n$ は連続であり，
  $0\leq f_n\leq1$ である．したがって $f_n\in C_b(\mathcal{X})$ である．

  $F$ が閉集合なので，$d(x,F)=0$ は $x\in F$ と同値である．
  よって $f_n(x)\to\mathbf{1}_F(x)$ がすべての $x\in \mathcal{X}$ で成り立つ．
  有界収束定理（定理~\ref{thm:meas-bounded-convergence}）と仮定より
  \[
    \eta(F)=\lim_{n\to\infty}\int f_n\,\d\eta
    =\lim_{n\to\infty}\int f_n\,\d\rho
    =\rho(F).
  \]
  したがって $\eta,\rho$ はすべての閉集合上で一致する．
  閉集合全体は $\Borel(\mathcal{X})$ を生成する $\pi$-系であり，
  $\eta(\mathcal{X})=\rho(\mathcal{X})=1$ だから，一致定理
  （定理~\ref{thm:meas-uniqueness}）より $\eta=\rho$ である．
\end{proof}

\begin{definition}{tight な確率測度族（Villani p.~43）}{meas-tight}
  Polish 空間 $\mathcal{X}$ 上の Borel 確率測度の族
  $\mathcal{F}\subset\Prob(\mathcal{X})$ が \textbf{tight} であるとは，
  任意の $\varepsilon>0$ に対し compact 集合
  $K_\varepsilon\subset\mathcal{X}$ が存在して
  \[
    \eta(\mathcal{X}\setminus K_\varepsilon)\leq\varepsilon
    \qquad(\forall\eta\in\mathcal{F})
  \]
  が成り立つことをいう．
\end{definition}

\begin{theorem}{Prokhorov の定理}{meas-prokhorov}
  $\mathcal{X}$ を Polish 空間とし，$\mathcal{F}\subset\Prob(\mathcal{X})$ を
  Borel 確率測度の族とする．このとき，次は同値である．
  \begin{enumerate}
    \item 任意の列 $(\eta_k)_{k\in\N}\subset\mathcal{F}$ に対し，
      ある部分列 $(\eta_{k_j})_{j\in\N}$ と
      $\eta\in\Prob(\mathcal{X})$ が存在して
      \[
        \eta_{k_j}\weakto\eta
      \]
      が成り立つ
      （すなわち，$\mathcal{F}$ は弱位相で precompact である）．
    \item $\mathcal{F}$ は tight である．
  \end{enumerate}
  （証明は省略する．Villani (2009) Chapter 4 で
  同じ形で引用されている古典的な定理であり，証明は
  Billingsley \emph{Convergence of Probability Measures}
  (2nd ed., 1999) Theorems 5.1, 5.2 にある．）
\end{theorem}

\section{$L^p$ ノルムと積分不等式}

\begin{definition}{$L^p$ ノルム}{meas-lp-norm}
  測度空間 $(\Omega,\mathcal{A},\mu)$ 上の可測関数 $f\colon\Omega\to\R$ と
  $1\leq p<\infty$ に対し
  \[
    \norm{f}_{L^p(\mu)}
    :=
    \left(\int_\Omega\abs{f}^p\,\d\mu\right)^{1/p}
  \]
  と定める．
\end{definition}

\begin{remark}{記法 $\norm{f}_{L^p(\mu)}$}{meas-lp-norm-notation}
  添字 $L^p(\mu)$ は，冪 $p$ と積分に用いる測度 $\mu$ の指定である：
  \[
    \norm{f}_{L^2(\mu)}=\left(\int_\Omega\abs{f}^2\,\d\mu\right)^{1/2},
    \qquad
    \norm{f}_{L^1(\nu)}=\int_\Omega\abs{f}\,\d\nu .
  \]
  同じ $f$ でも測度が異なれば値は異なる．
  本編では coupling $\pi\in\Couplings(\mu,\nu)$ に対する
  \[
    \norm{d}_{L^p(\pi)}
    =\left(\int_{\mathcal{X}\times \mathcal{X}}d(x,y)^p\,\d\pi(x,y)\right)^{1/p}
  \]
  の形で用い，
  \[
    \Wass_p(\mu,\nu)=\inf_{\pi\in\Couplings(\mu,\nu)}\norm{d}_{L^p(\pi)}
  \]
  と書ける（注意~\ref{rem:wp-welldef}）．
\end{remark}

\begin{theorem}{Hölder の不等式}{meas-holder}
  $1<r,r'<\infty$，$1/r+1/r'=1$ とする．可測関数 $f,g$ に対し
  \[
    \int_\Omega\abs{fg}\,\d\mu
    \leq
    \norm{f}_{L^r(\mu)}\,\norm{g}_{L^{r'}(\mu)}.
  \]
  （標準的な結果．証明は省略する．Dudley (2002) §5.1 参照．）
\end{theorem}

\begin{claim}{$L^p$ ノルムの単調性}{meas-lp-monotone}
  可測関数 $f,g$ が $\abs{f}\leq\abs{g}$ $\mu$-a.e. をみたすならば，
  $1\leq p<\infty$ に対し
  $\norm{f}_{L^p(\mu)}\leq\norm{g}_{L^p(\mu)}$ である．
\end{claim}

\begin{proof}
  $\abs{f}^p\leq\abs{g}^p$ $\mu$-a.e. だから，
  積分の単調性より
  $\int\abs{f}^p\,\d\mu\leq\int\abs{g}^p\,\d\mu$．
  両辺の $1/p$ 乗を取ればよい．
\end{proof}

\begin{theorem}{Minkowski の不等式}{meas-minkowski}
  可測関数 $f,g$ と $1\leq p<\infty$ に対し
  \[
    \left(\int_\Omega\abs{f+g}^p\,\d\mu\right)^{1/p}
    \leq
    \left(\int_\Omega\abs{f}^p\,\d\mu\right)^{1/p}
    +
    \left(\int_\Omega\abs{g}^p\,\d\mu\right)^{1/p}.
  \]
  （Hölder の不等式から従う標準的な結果．証明は省略する．
  例えば Dudley (2002) §5.1 参照．）
\end{theorem}

\section{像測度}

\begin{definition}{像測度}{meas-pushforward}
  可測空間 $(\Omega,\mathcal{A})$，$(\Omega',\mathcal{A}')$ と
  可測写像 $T\colon\Omega\to\Omega'$，
  $(\Omega,\mathcal{A})$ 上の測度 $\mu$ に対し，
  \[
    (\mu\circ T^{-1})(B):=\mu(T^{-1}(B))
    \qquad(B\in\mathcal{A}')
  \]
  と定める．$\mu\circ T^{-1}$ を $T$ による $\mu$ の\textbf{像測度}という．
  ここで $T^{-1}(B):=\{\omega\in\Omega\mid T(\omega)\in B\}$ は
  集合 $B$ の\textbf{逆像}であり，$T$ の可逆性は仮定しない．
  逆像は互いに素な可算合併を互いに素な可算合併に写すから，
  $\mu\circ T^{-1}$ の $\sigma$-加法性は $\mu$ のそれから従い，
  $\mu\circ T^{-1}$ は $(\Omega',\mathcal{A}')$ 上の測度である．
  $\mu$ が確率測度なら
  $(\mu\circ T^{-1})(\Omega')=\mu(\Omega)=1$ より
  $\mu\circ T^{-1}$ も確率測度である．
  なお文献では像測度を $T_\#\mu$ や $T_*\mu$ とも書き，
  押し出し測度ともよぶ（Villani 2009，桑江ほか 2015 は $T_\#\mu$）．
\end{definition}

\begin{lemma}{変数変換公式}{meas-change-of-variables}
  $T\colon\Omega\to\Omega'$ を可測写像，$\mu$ を $\Omega$ 上の測度とする．
  非負可測関数 $f\colon\Omega'\to[0,+\infty]$ に対し
  \[
    \int_{\Omega'}f\,\d(\mu\circ T^{-1})
    =
    \int_\Omega f\circ T\,\d\mu
  \]
  が成り立つ．また可測な $f\colon\Omega'\to\R$ については，
  一方の積分が絶対収束すれば他方もそうであり，同じ等式が成り立つ．
\end{lemma}

\begin{proof}
  $f=\mathbf{1}_B$（$B\in\mathcal{A}'$）のとき：
  $\mathbf{1}_B\circ T=\mathbf{1}_{T^{-1}(B)}$ だから
  \[
    \int_{\Omega'}\mathbf{1}_B\,\d(\mu\circ T^{-1})
    =(\mu\circ T^{-1})(B)
    =\mu(T^{-1}(B))
    =\int_\Omega\mathbf{1}_B\circ T\,\d\mu.
  \]
  $f$ が非負単関数のとき：指示関数の非負線形結合だから，
  積分の線形性より従う．
  $f$ が非負可測のとき：非負単関数列 $s_n\uparrow f$ を取れば
  （注意~\ref{rem:meas-measurable-ops}），
  $s_n\circ T\uparrow f\circ T$（各点）であり，
  単調収束定理（定理~\ref{thm:meas-mct}）を両辺に適用して
  \[
    \int f\,\d(\mu\circ T^{-1})
    =\lim_n\int s_n\,\d(\mu\circ T^{-1})
    =\lim_n\int s_n\circ T\,\d\mu
    =\int f\circ T\,\d\mu.
  \]
  一般の可測な $f$ のとき：$f=f^+-f^-$（正部分・負部分）と分解すると，
  非負の場合から
  $\int f^\pm\,\d(\mu\circ T^{-1})=\int f^\pm\circ T\,\d\mu$ であり，
  絶対収束の条件も両辺で一致する．差を取れば等式を得る．
\end{proof}

$\mathcal{X},\mathcal{Y}$ を可分距離空間とする．
直積 $\mathcal{X}\times \mathcal{Y}$ には距離
$d_{\mathcal{X}\times \mathcal{Y}}((x,y),(x',y')):=\max\{d_X(x,x'),d_Y(y,y')\}$
を入れる（稠密可算集合の直積が稠密なので，$\mathcal{X}\times \mathcal{Y}$ も可分である）．
また $\mathcal{X},\mathcal{Y}$ が完備なら $\mathcal{X}\times \mathcal{Y}$ も完備であるから，
$\mathcal{X},\mathcal{Y}$ が Polish 空間なら $\mathcal{X}\times \mathcal{Y}$ も Polish 空間である．
3個以上の有限個の直積 $\mathcal{X}_1\times\cdots\times\mathcal{X}_k$ にも
同様に $\max$ 距離を入れ，可分性・完備性は保たれる
（2因子の場合の繰り返し）．

\begin{remark}{直積空間の Borel $\sigma$-代数（前提とする事実）}{meas-borel-product}
  $\mathcal{X},\mathcal{Y}$ が可分距離空間のとき，
  $\Borel(\mathcal{X}\times \mathcal{Y})$ は矩形集合
  $\{A\times B\mid A\in\Borel(\mathcal{X}),\,B\in\Borel(\mathcal{Y})\}$
  の生成する $\sigma$-代数
  $\Borel(\mathcal{X})\otimes\Borel(\mathcal{Y})$ と一致する（証明は省略する）．
  なお矩形集合の族は
  $(A\times B)\cap(A'\times B')=(A\cap A')\times(B\cap B')$
  より $\pi$-系である．
  有限個の直積 $\mathcal{X}_1\times\cdots\times\mathcal{X}_k$ についても同様に，
  矩形集合 $A_1\times\cdots\times A_k$ の全体は $\pi$-系をなし，
  $\Borel(\mathcal{X}_1\times\cdots\times\mathcal{X}_k)$ を生成する
  （2因子の場合を繰り返せばよい）．
\end{remark}

\begin{lemma}{周辺積分公式}{meas-marginal-integral}
  $\mathcal{X},\mathcal{Y}$ を距離空間とし，
  $\pi\in\Prob(\mathcal{X}\times\mathcal{Y})$ の第1周辺分布を
  $\mu\in\Prob(\mathcal{X})$ とする．すなわち，
  \[
    \pi(A\times\mathcal{Y})=\mu(A)
    \qquad(\forall A\in\Borel(\mathcal{X}))
  \]
  とする．このとき，任意の非負可測関数
  $f\colon\mathcal{X}\to[0,+\infty]$，または任意の $\mu$-可積分関数
  $f\colon\mathcal{X}\to\R$ に対して，
  \[
    \int_{\mathcal{X}\times \mathcal{Y}}f(x)\,\d\pi(x,y)
    =
    \int_{\mathcal{X}}f(x)\,\d\mu(x)
  \]
  が成り立つ．第2周辺分布についても同様である．
\end{lemma}

\begin{proof}
  座標写像 $(x,y)\mapsto x$ は連続だから，$f(x)$ は可測である．
  $f=\mathbf{1}_A$（$A\in\Borel(\mathcal{X})$）ならば，
  \[
    \int_{\mathcal{X}\times\mathcal{Y}}\mathbf{1}_A(x)\,\d\pi(x,y)
    =\pi(A\times\mathcal{Y})
    =\mu(A)
    =\int_{\mathcal{X}}\mathbf{1}_A\,\d\mu.
  \]
  したがって，積分の線形性と単調収束定理
  （定理~\ref{thm:meas-mct}）により
  \[
    \mathbf{1}_A
    \quad\Longrightarrow\quad
    \text{非負単関数}
    \quad\Longrightarrow\quad
    \text{非負可測関数}
  \]
  の順に積分公式が従う．$f$ が $\mu$-可積分ならば，
  $f=f^+-f^-$ と分解して非負可測関数の場合を適用すればよい．
\end{proof}

\section{Lebesgue 測度}

\begin{definition}{Euclid 空間上の Lebesgue 測度}{meas-lebesgue}
  $n\in\N$ とする．半開矩形の体積を
  \[
    \lambda_n\left(\prod_{i=1}^n(a_i,b_i]\right)
    :=\prod_{i=1}^n(b_i-a_i)
    \qquad(a_i<b_i)
  \]
  と定め，これを $\Borel(\R^n)$ 上へ拡張して得られる測度
  $\lambda_n$ を $\R^n$ 上の\textbf{Lebesgue 測度}という．
  その存在と一意性は測度論の標準定理として認める．
  $\lambda_n$ に関する積分を
  \[
    \int_{\R^n}f(x)\,\d x
    :=\int_{\R^n}f\,\d\lambda_n
  \]
  と略記する．連続関数の（広義）Riemann 積分と Lebesgue 積分の一致，
  および微積分学の基本定理も既知として用いる．
\end{definition}

\section{直積測度と測度の分解}

\begin{theorem}{直積確率測度と Tonelli の定理}{meas-product}
  $\mathcal{X},\mathcal{Y}$ を可分距離空間，
  $\mu\in\Prob(\mathcal{X})$，$\nu\in\Prob(\mathcal{Y})$ とする．
  $\Borel(\mathcal{X}\times \mathcal{Y})$ 上の確率測度 $\mu\otimes\nu$ で
  \[
    (\mu\otimes\nu)(A\times B)=\mu(A)\,\nu(B)
    \qquad(A\in\Borel(\mathcal{X}),\;B\in\Borel(\mathcal{Y}))
  \]
  をみたすものがただ一つ存在する（\textbf{直積測度}）．
  さらに非負可測関数 $f\colon \mathcal{X}\times \mathcal{Y}\to[0,\infty)$ に対し
  \[
    \int_{\mathcal{X}\times \mathcal{Y}}f\,\d(\mu\otimes\nu)
    =\int_{\mathcal{X}}\!\!\int_{\mathcal{Y}} f(x,y)\,\d\nu(y)\,\d\mu(x)
    =\int_{\mathcal{Y}}\!\!\int_{\mathcal{X}} f(x,y)\,\d\mu(x)\,\d\nu(y)
  \]
  が成り立つ（内側の積分は可測関数を定める）．
  有限個の確率測度の直積
  $\mu_1\otimes\cdots\otimes\mu_n$ についても同様である．
  （標準的な結果．証明は省略する．Dudley (2002) §4.4 参照．）
\end{theorem}

\begin{theorem}{測度分解定理（正則条件付き確率）}{meas-disintegration}
  $\mathcal{X},\mathcal{Y}$ を Polish 空間，$\pi\in\Prob(\mathcal{X}\times \mathcal{Y})$ とし，
  $\nu$ を $\pi$ の第2周辺分布とする．
  このとき写像 $K\colon \mathcal{Y}\times\Borel(\mathcal{X})\to[0,1]$ で
  次をみたすものが存在する：
  \begin{enumerate}
    \item 各 $y\in \mathcal{Y}$ で $K(y,\cdot)\in\Prob(\mathcal{X})$，
    \item 各 $A\in\Borel(\mathcal{X})$ で $y\mapsto K(y,A)$ は可測，
    \item すべての $A\in\Borel(\mathcal{X})$，$B\in\Borel(\mathcal{Y})$ に対し
      \[
        \pi(A\times B)=\int_B K(y,A)\,\d\nu(y).
      \]
  \end{enumerate}
  (i)(ii) をみたす $K$ を（$\mathcal{Y}$ から $\mathcal{X}$ への）\textbf{Markov 核}といい，
  (iii) を $\pi(\d x,\d y)=K(y,\d x)\,\nu(\d y)$ と書く．
  $K$ は $\nu$-a.e. の意味で一意である．
  （Polish 空間上の標準的な結果．証明は省略する．
  Dudley (2002) Chapter 10，Villani (2009) Chapter 1，
  桑江ほか (2015) 補題 1.32 参照．）
\end{theorem}

\begin{theorem}{Markov 核による測度の貼り合わせ（前提とする事実）}{meas-kernel-composition}
  $\mathcal{X},\mathcal{Y},\mathcal{Z}$ を Polish 空間，$\nu\in\Prob(\mathcal{Y})$ とし，
  $K_1$ を $\mathcal{Y}$ から $\mathcal{X}$ への，$K_2$ を $\mathcal{Y}$ から $\mathcal{Z}$ への Markov 核とする．
  このとき $\Borel(\mathcal{X}\times \mathcal{Y}\times \mathcal{Z})$ 上の確率測度 $\gamma$ で
  \[
    \gamma(A\times B\times C)
    =\int_B K_1(y,A)\,K_2(y,C)\,\d\nu(y)
    \qquad(A\in\Borel(\mathcal{X}),\,B\in\Borel(\mathcal{Y}),\,C\in\Borel(\mathcal{Z}))
  \]
  をみたすものがただ一つ存在する．
  これを $\gamma(\d x,\d y,\d z)=K_1(y,\d x)\,\nu(\d y)\,K_2(y,\d z)$ と書く．
  （直積測度の構成と同様の標準的な手順（単調族論法）による．
  証明は省略する．一意性は矩形集合が $\pi$-系をなすことと
  一致定理（定理~\ref{thm:meas-uniqueness}）から従う．）
\end{theorem}
